\sgasection{3}{\texorpdfstring{The category of $\mathbf{O}$-modules, the category of $\mathbf{G}$-$\mathbf{O}$-modules}{The category of O-modules, the category of G-O-modules}}
\label{I.sec.3}

\begin{definition}{3.1}\label{I.3.1}
Let \(\mathbf{O}\) and \(\mathbf{F}\) be two objects of \(\widehat{\mathcal{C}}\).
One says that \(\mathbf{F}\) is a \(\widehat{\mathcal{C}}\)-module over the
\(\widehat{\mathcal{C}}\)-ring \(\mathbf{O}\), or for short an \(\mathbf{O}\)-module,
if, for every \(S\in\Ob(\mathcal{C})\), one has endowed \(\mathbf{O}(S)\) with a
ring structure and \(\mathbf{F}(S)\) with a module structure over this ring, in
such a way that for every arrow \(S'\to S''\) of \(\mathcal{C}\),
\(\mathbf{O}(S'')\to\mathbf{O}(S')\) is a ring homomorphism and
\(\mathbf{F}(S'')\to\mathbf{F}(S')\) a homomorphism of abelian groups compatible
with this ring homomorphism.

If \(\mathbf{O}\) is fixed, one defines in the usual way a morphism of the
\(\mathbf{O}\)-modules \(\mathbf{F}\) and \(\mathbf{F}'\), whence the
commutative group \(\Hom_{\mathbf{O}}(\mathbf{F},\mathbf{F}')\), and the
category of \(\mathbf{O}\)-modules, denoted \((\mathbf{O}\text{-Mod.})\).
\end{definition}

\begin{lemma}{3.1.1}\label{I.3.1.1}\ndemark
\((\mathbf{O}\text{-Mod.})\) is endowed with the structure of an abelian
category, defined \og argument by argument\fg. Moreover,
\((\mathbf{O}\text{-Mod.})\) satisfies axiom~(AB~5) (\cf [Gr57, 1.5]), \ie
arbitrary direct sums exist and if \(\mathbf{M}\) is an \(\mathbf{O}\)-module,
\(\mathbf{N}\) a submodule, and \((\mathbf{F}_i)_{i\in I}\) an increasing
filtered family of submodules of \(\mathbf{M}\), then
\[
\bigcup_{i\in I}(\mathbf{F}_i\cap\mathbf{N})
=\Bigl(\bigcup_{i\in I}\mathbf{F}_i\Bigr)\cap\mathbf{N}.
\]
\end{lemma}
\ndetext{Statements 3.1.1 and 3.1.2 have been added in order to show that
the category \((\mathbf{O}\text{-Mod.})\) is abelian and satisfies
axiom~(AB~5), and moreover has enough injective objects if the category
\(\mathcal{C}\) is small. An error in 3.1.2, pointed out in 2016 by Linyuan
Liu, is corrected here.}

Indeed, let \(f\colon\mathbf{F}\to\mathbf{F}'\) be a morphism of
\(\mathbf{O}\)-modules. One defines the \(\mathbf{O}\)-modules \(\Ker f\)
(\resp\ \(\Im f\) and \(\Coker f\)) by setting, for every
\(S\in\Ob(\mathcal{C})\), \((\Ker f)(S)=\Ker f(S)\) (\resp\ \(\cdots\)).
Then \(\Ker f\) (\resp\ \(\Coker f\)) is a kernel (\resp\ a cokernel) of
\(f\), and one has an isomorphism of \(\mathbf{O}\)-modules
\(\mathbf{F}/\Ker f\isomto\Im f\). This proves that
\((\mathbf{O}\text{-Mod.})\) is an abelian category.

Arbitrary direct sums exist and are defined \og argument by argument\fg.
Finally, if \(\mathbf{M}\) is an \(\mathbf{O}\)-module, \(\mathbf{N}\) a
submodule, and \((\mathbf{F}_i)_{i\in I}\) an \emph{increasing filtered
family} of submodules of \(\mathbf{M}\), then the inclusion
\[
\bigcup_{i\in I}(\mathbf{F}_i\cap\mathbf{N})
\subset\Bigl(\bigcup_{i\in I}\mathbf{F}_i\Bigr)\cap\mathbf{N}
\]
is an equality: indeed, if \(S\in\Ob(\mathcal{C})\) and
\(x\in\mathbf{N}(S)\cap\bigcup_i\mathbf{F}_i(S)\), there exists \(i\in I\)
such that \(x\in\mathbf{N}(S)\cap\mathbf{F}_i(S)\).

\begin{proposition}{3.1.2}\label{I.3.1.2}\footnotemark[\value{footnote}]
One assumes that the category \(\mathcal{C}\) is small, \ie that
\(\Ob(\mathcal{C})\) is a set. Then \((\mathbf{O}\text{-Mod.})\) has a
generator, and hence enough injective objects.
\end{proposition}

\begin{proof}
For every object \(S\) of \(\mathcal{C}\), one defines the
\(\mathbf{O}\)-module \(\mathbf{O}_S\) as follows.
For every object \(S'\), \(\mathbf{O}_S(S')\) is a direct sum of copies of
\(\mathbf{O}(S')\) indexed by \(\Hom(S',S)\), \ie
\[
\mathbf{O}_S(S')=\bigoplus_{g\colon S'\to S}\mathbf{O}(S')\,g.
\]
For every morphism \(f\colon S''\to S'\), denoting by \(f^{*}\) the ring
homomorphism \(\mathbf{O}(S')\to\mathbf{O}(S'')\), the morphism
\(\mathbf{O}_S(f)\colon\mathbf{O}_S(S')\to\mathbf{O}_S(S'')\) sends each
element \(ag\) of \(\mathbf{O}_S(S')\) (where \(a\in\mathbf{O}(S')\) and
\(g\colon S'\to S\)) to the element \(f^{*}(a)\,g\circ f\) of
\(\mathbf{O}_S(S'')\). One easily checks that this does indeed define a
functor in \(\mathbf{O}\)-modules.

Since, by hypothesis, \(\Ob(\mathcal{C})\) is a set, one may consider the
direct sum \(\mathbf{G}=\bigoplus_{S\in\Ob(\mathcal{C})}\mathbf{O}_S\).

\begin{lemma}{3.1.2.1}\label{I.3.1.2.1}
Let \(\mathbf{F}\) be an \(\mathbf{O}\)-module. For every
\(S\in\Ob(\mathcal{C})\), every element \(x\) of \(\mathbf{F}(S)\) defines a
unique morphism of \(\mathbf{O}\)-modules \(x\colon\mathbf{O}_S\to\mathbf{F}\)
such that \(x(1\,\mathrm{id}_S)=x\).
\end{lemma}

\begin{proof}
Let \(\phi\) be a morphism of \(\mathbf{O}\)-modules
\(\mathbf{O}_S\to\mathbf{F}\) such that \(\phi(1\,\mathrm{id}_S)=x\).
Let \(g\colon S'\to S\) and \(a\in\mathbf{O}(S')\).
The \(\mathbf{O}(S')\)-linearity and the commutativity of the diagram
\[
\xymatrix{
\mathbf{O}_S(S) \ar[r]^{\phi} \ar[d]_{\mathbf{O}_S(g)}
  & \mathbf{F}(S) \ar[d]^{\mathbf{F}(g)} \\
\mathbf{O}_S(S') \ar[r]^{\phi}
  & \mathbf{F}(S')
}
\]
imply that
\(\phi(ag)=a\phi(1g)=a\mathbf{F}(g)(\phi(1\,\mathrm{id}_S))=a\mathbf{F}(g)(x)\).
Thus \(\phi\), if it exists, is determined by the preceding equality.
Conversely, if one defines \(\phi\) in this way, then for every
\(f\colon S''\to S'\) one has
\[
\phi\circ\mathbf{O}_S(f)(ag)
=\phi(f^{*}(a)\,g\circ f)
=f^{*}(a)\mathbf{F}(g\circ f)(x)
=\mathbf{F}(f)(a\mathbf{F}(g)(x))
=\mathbf{F}(f)\circ\phi(ag),
\]
\ie the diagram
\[
\xymatrix{
\mathbf{O}_S(S') \ar[r]^{\phi} \ar[d]_{\mathbf{O}_S(f)}
  & \mathbf{F}(S') \ar[d]^{\mathbf{F}(f)} \\
\mathbf{O}_S(S'') \ar[r]^{\phi}
  & \mathbf{F}(S'')
}
\]
is commutative. This proves the lemma.
\end{proof}

Now let \(\mathbf{F}\) be an \(\mathbf{O}\)-module. For every
\(S\in\Ob(\mathcal{C})\), let \(F_0(S)\) be a system of generators of the
\(\mathbf{O}(S)\)-module \(\mathbf{F}(S)\), and let \(\mathbf{L}_S\) be the
direct sum, indexed by \(F_0(S)\), of copies of \(\mathbf{O}_S\). By the
lemma, one obtains a morphism of \(\mathbf{O}\)-modules
\(\mathbf{L}_S\to\mathbf{F}\) which is such that the morphism
\(\mathbf{L}_S(S)\to\mathbf{F}(S)\) is surjective, hence an epimorphism of
the category of \(\mathbf{O}(S)\)-modules.

Set \(I=\bigsqcup_{S\in\Ob(\mathcal{C})}F_0(S)\). Then one obtains a morphism
of \(\mathbf{O}\)-modules
\[
\mathbf{G}^{\oplus I}
\longrightarrow\bigoplus_{S\in\Ob(\mathcal{C})}\mathbf{L}_S
\longrightarrow\mathbf{F}
\]
which is an epimorphism \og argument by argument\fg, hence an epimorphism
of \((\mathbf{O}\text{-Mod.})\). This shows that \(\mathbf{G}\) is a
generator of \((\mathbf{O}\text{-Mod.})\) (\cf [Gr57, 1.9.1]). Since
\((\mathbf{O}\text{-Mod.})\) satisfies (AB~5), it then follows from
[Gr57, 1.10.1] that \((\mathbf{O}\text{-Mod.})\) has enough injective
objects.
\end{proof}

\begin{remark}{3.1.3}\label{I.3.1.3}
If \(\mathbf{O}_0\) is the \(\widehat{\mathcal{C}}\)-ring defined by
\(\mathbf{O}_0(S)=\mathbb{Z}\) (which must not be confused with the functor
associated with the constant object \(\mathbb{Z}\)), then the category of
\(\mathbf{O}_0\)-modules is isomorphic to the category of commutative
\(\widehat{\mathcal{C}}\)-groups.
\end{remark}

\begin{definition}{3.1.4}\label{I.3.1.4}
Let us observe that, if \(\mathbf{F}\) is an \(\mathbf{O}\)-module, then for
every \(S\in\Ob(\mathcal{C})\), \(\mathbf{F}_S\) is an
\(\mathbf{O}_S\)-module. One may therefore define an abelian
\(\widehat{\mathcal{C}}\)-group \(\Hom_{\mathbf{O}}(\mathbf{F},\mathbf{F}')\)
by
\[
\Hom_{\mathbf{O}}(\mathbf{F},\mathbf{F}')(S)
=\Hom_{\mathbf{O}_S}(\mathbf{F}_S,\mathbf{F}'_S).
\]
One will likewise define objects
\[
\Isom_{\mathbf{O}}(\mathbf{F},\mathbf{F}'),
\qquad
\End_{\mathbf{O}}(\mathbf{F})
\quad\text{and}\quad
\Aut_{\mathbf{O}}(\mathbf{F}),
\]
the last being a \(\widehat{\mathcal{C}}\)-group for the structure induced
by the composition of automorphisms.
\end{definition}

\begin{definition}{3.2}\label{I.3.2}
Let \(\mathbf{O}\) be a \(\widehat{\mathcal{C}}\)-ring, \(\mathbf{F}\) an
\(\mathbf{O}\)-module, and \(\mathbf{G}\) a \(\widehat{\mathcal{C}}\)-group.
A structure of \(\mathbf{G}\)-object such that, for every
\(S\in\Ob(\mathcal{C})\) and every \(g\in\mathbf{G}(S)\), the automorphism of
\(\mathbf{F}(S)\) defined by \(g\) is an automorphism of its structure of
\(\mathbf{O}(S)\)-module, is called a \emph{structure of
\(\mathbf{G}\)-\(\mathbf{O}\)-module on \(\mathbf{F}\)}.

It amounts to the same to say that the morphism of
\(\widehat{\mathcal{C}}\)-groups
\[
\rho\colon\mathbf{G}\longrightarrow\Aut(\mathbf{F})
\]
defined in~2.3 sends \(\mathbf{G}\) into the
\(\widehat{\mathcal{C}}\)-subgroup \(\Aut_{\mathbf{O}}(\mathbf{F})\) of
\(\Aut(\mathbf{F})\). To give a \emph{structure of
\(\mathbf{G}\)-\(\mathbf{O}\)-module} on the \(\mathbf{O}\)-module
\(\mathbf{F}\) is thus to give a \emph{morphism of
\(\widehat{\mathcal{C}}\)-groups}
\[
\rho\colon\mathbf{G}\longrightarrow\Aut_{\mathbf{O}}(\mathbf{F}).
\]

One defines in a natural way the abelian group
\(\Hom_{\mathbf{G}\text{-}\mathbf{O}}(\mathbf{F},\mathbf{F}')\), hence the
additive category of \(\mathbf{G}\)-\(\mathbf{O}\)-modules, denoted
\((\mathbf{G}\text{-}\mathbf{O}\text{-Mod.})\).
\end{definition}

\begin{remark}{3.2.1}\label{I.3.2.1}
The reader will observe that \((\mathbf{G}\text{-}\mathbf{O}\text{-Mod.})\)
may equally be defined as the category of
\(\mathbf{O}[\mathbf{G}]\)-modules, where \(\mathbf{O}[\mathbf{G}]\) is the
algebra of the \(\widehat{\mathcal{C}}\)-group \(\mathbf{G}\) over the
\(\widehat{\mathcal{C}}\)-ring \(\mathbf{O}\), whose definition is
clear.\nde{The following sentence has been added; this will be used in
Section~5.}
Consequently, by~3.1.1 and~3.1.2,
\((\mathbf{G}\text{-}\mathbf{O}\text{-Mod.})\) is an abelian category
satisfying axiom~(AB~5); moreover, if \(\mathcal{C}\) is \emph{small},
\((\mathbf{G}\text{-}\mathbf{O}\text{-Mod.})\) has enough injective objects.
\end{remark}

All the constructions of this paragraph specialize at once to the case where
\(\mathbf{G}\) (or \(\mathbf{O}\), or both) are representable by objects of
\(\mathcal{C}\) which are thereby endowed with the corresponding algebraic
structures.

We have treated succinctly the case of the principal algebraic structures
encountered in the rest of this seminar. For the others (the structure of
\(\mathbf{O}\)-Lie algebra, for example), we believe that the reader will
have had enough examples in this paragraph to be able, in each particular
case, to put into operation for himself the general mechanism sketched
in~2.2.

We are now going to apply what we have just done to the category of schemes,
denoted~\(\Sch\), and more generally to the categories \(\Sch_{/S}\) (which
one will also denote \((\textup{Sch}/S)\)).

\sgasection{4}{Algebraic structures in the category of schemes}
\label{I.sec.4}

We shall permit ourselves, whenever there is no ambiguity, the following
abuses of language: one will denote by \(T\) the object
\(T\xrightarrow{f}S\) of \(\mathcal{C}_{/S}\), the datum of \(f\) (\og
structural morphism of \(T\)\fg) being understood, and one will identify
\(\mathcal{C}\) with a subcategory of \(\widehat{\mathcal{C}}\). In view of
the compatibilities stated in the preceding paragraphs, these
identifications may be made without danger.

On the other hand, we shall simplify the terminology on the following model:
a \(\Sch\)-group will also be called a \emph{group scheme}, a
\(\Sch_{/S}\)-group a \emph{group scheme over \(S\)}, or an \emph{\(S\)-group
scheme}, or an \emph{\(S\)-group}, or an \emph{\(A\)-group} when \(S\) is the
spectrum of the ring \(A\).

\par\addvspace{0.55\baselineskip}
\noindent\textbf{4.1. Constant schemes.}\label{I.4.1}\ ---
The category of schemes satisfies the conditions required in~1.8. One may
therefore define constant objects there; for every set \(E\), one has a
scheme \(E_{\mathbb{Z}}\) and, for every scheme \(S\), an \(S\)-scheme
\(E_S=(E_{\mathbb{Z}})_S\). Let us recall (\cf\loccit) that for every
\(S\)-scheme \(T\), \(\Hom_S(T,E_S)\) is the set of locally constant maps from
the topological space \(T\) to \(E\).

The functor \(E\mapsto E_S\) commutes with finite projective limits. In
particular, if \(G\) is a group, \(G_S\) is an \emph{\(S\)-group scheme}; if
\(O\) is a ring, \(O_S\) is an \emph{\(S\)-ring scheme}, etc.

\par\addvspace{0.55\baselineskip}
\noindent\textbf{4.2. Affine \(S\)-groups.}\label{I.4.2}\ ---
Let us recall a certain number of things about affine \(S\)-schemes
(EGA~II, \S~1). One says that the \(S\)-scheme \(T\) is \emph{affine over
\(S\)} if the inverse image of every affine open of \(S\) is affine. The
\(\OS\)-algebra \(f_*(\mathcal{O}_T)\), which one denotes by \(\mathcal{A}(T)\),
is then \emph{quasi-coherent} (\(f\) denotes the structural morphism of
\(T\)). Conversely, to every quasi-coherent \(\OS\)-algebra \(\mathcal{A}\),
one may associate an \(S\)-scheme affine over \(S\), denoted
\(\Spec(\mathcal{A})\). These functors \(T\mapsto\mathcal{A}(T)\) and
\(\mathcal{A}\mapsto\Spec(\mathcal{A})\) are equivalences quasi-inverse to
one another between the category of \(S\)-schemes affine over \(S\) and the
category opposite to that of quasi-coherent \(\OS\)-algebras.

It follows that to give an algebraic structure on an \(S\)-scheme \(T\)
affine over \(S\) is equivalent to giving the corresponding structure on
\(\mathcal{A}(T)\) in the category opposite to that of quasi-coherent
\(\OS\)-algebras. In particular, if \(G\) is an affine \(S\)-group over \(S\),
\(\mathcal{A}(G)\) is endowed with the structure of an augmented
\(\OS\)-bialgebra, that is, one has two \emph{morphisms of
\(\OS\)-algebras}
\begin{gather*}
\Delta\colon\mathcal{A}(G)\longrightarrow
  \mathcal{A}(G)\otimes_{\OS}\mathcal{A}(G)
\qquad\text{and}\qquad
\varepsilon\colon\mathcal{A}(G)\longrightarrow\OS
\end{gather*}
corresponding to the morphisms of \(S\)-schemes
\[
\pi\colon G\times G\longrightarrow G
\qquad\text{and}\qquad
e_G\colon S\longrightarrow G.
\]
The maps \(\Delta\) and \(\varepsilon\) satisfy the following conditions
(which express that \(G\) is an \(S\)-monoid):\nde{And, of course, the
inversion morphism \(G\to G\) induces a morphism of \(\OS\)-algebras
\(\tau\colon\mathcal{A}(G)\to\mathcal{A}(G)\) which, together with \(\Delta\)
and \(\varepsilon\), makes \(\mathcal{A}(G)\) into a Hopf \(\OS\)-algebra.}

(HA~1) \(\Delta\) \emph{is coassociative}: the following diagram is commutative
\[
\xymatrix@C=3.2em@R=2.6em{
& \mathcal{A}(G)\otimes_{\OS}\mathcal{A}(G)
    \ar[dr]^{\mathrm{Id}\otimes\Delta} & \\
\mathcal{A}(G) \ar[ur]^{\Delta} \ar[dr]_{\Delta}
&& \mathcal{A}(G)\otimes_{\OS}\mathcal{A}(G)\otimes_{\OS}\mathcal{A}(G) \\
& \mathcal{A}(G)\otimes_{\OS}\mathcal{A}(G)
    \ar[ur]_{\Delta\otimes\mathrm{Id}} &
}\,.
\]

(HA~2): \(\Delta\) \emph{is compatible with \(\varepsilon\)}: the following
two composites are the identity
\begin{gather*}
\mathcal{A}(G)\xrightarrow{\Delta}
  \mathcal{A}(G)\otimes_{\OS}\mathcal{A}(G)
  \xrightarrow{\mathrm{Id}\otimes\varepsilon}
  \mathcal{A}(G)\otimes_{\OS}\OS
  \isomto\mathcal{A}(G)\,,\\
\mathcal{A}(G)\xrightarrow{\Delta}
  \mathcal{A}(G)\otimes_{\OS}\mathcal{A}(G)
  \xrightarrow{\varepsilon\otimes\mathrm{Id}}
  \OS\otimes_{\OS}\mathcal{A}(G)
  \isomto\mathcal{A}(G)\,.
\end{gather*}

Let us take the opportunity to remark once again that it follows from the
definition of an \(S\)-group structure that, in order to give such a
structure on an \(S\)-scheme \(G\) affine over \(S\), it is not necessary to
check anything at all on \(\mathcal{A}(G)\), but simply to endow each
\(G(S')\), for \(S'\) over \(S\), with a group structure functorial in
\(S'\).
% typo?: kept as printed "quoi ce soit"
This remark applies \emph{mutatis mutandis} to morphisms.

\par\addvspace{0.7\baselineskip}
\noindent\textbf{4.3. The groups \(\mathbb{G}_a\) and \(\mathbb{G}_m\). The ring \(\mathbb{O}\)}\label{I.4.3}

\numpara{4.3.1}\label{I.4.3.1}
Let \(\Ga\) be the functor from \(\Sch^{\circ}\) to \(\Ens\) defined by
\[
\Ga(S)=\Gamma(S,\OS),
\]
endowed with the structure of \((\widehat{\textup{Sch}})\)-group defined by
the additive group structure of the ring \(\Gamma(S,\OS)\). It is
representable by an affine scheme which one will denote by \(\mathbb{G}_a\),
and which is therefore a \emph{group scheme}
\[
\mathbb{G}_a=\Spec\mathbb{Z}[T].
\]
Indeed,
\(\Ga(S)=\Hom(S,\mathbb{G}_a)
=\Hom_{\mathrm{Alg.}}(\mathbb{Z}[T],\Gamma(S,\OS))
\simeq\Gamma(S,\OS)\).

For every scheme \(S\), one therefore has an affine \(S\)-group over \(S\),
denoted \(\mathbb{G}_{a,S}\), which corresponds to the \(\OS\)-bialgebra
\(\OS[T]\), with the diagonal map and the augmentation defined by
\[
\Delta(T)=T\otimes 1+1\otimes T,
\qquad
\varepsilon(T)=0.
\]

\numpara{4.3.2}\label{I.4.3.2}
Let \(\Gm\) be the functor from \(\Sch^{\circ}\) to \(\Ens\) defined by
\[
\Gm(S)=\Gamma(S,\OS)^{\times},
\]
where \(\Gamma(S,\OS)^{\times}\) denotes the multiplicative group of
invertible elements of the ring \(\Gamma(S,\OS)\), endowed with its natural
structure of \((\widehat{\textup{Sch}})\)-group.
